% Section 4.2, from lectures/L04/appendix/b_factory_smart_pointers.md.

\section{Driver Factory with Smart Pointers (Abstract Factory)}
\label{c4:sec:smart}\label{c4:app:b}
In this section we further develop the factory example from \secref{c4:sec:raw}.

We replace raw pointers with smart pointers and use a simplified form of the \textbf{Abstract
Factory pattern}.

The purpose is to:
\begin{itemize}
  \item Avoid manual memory management.
  \item Clarify ownership.
  \item Make the system safer and more robust.
\end{itemize}

\subsection{What is Abstract Factory?}
\label{c4:sec:abstractfactory}
\textbf{Abstract Factory} is a design pattern where:
\begin{itemize}
  \item A common factory interface defines how objects are created.
  \item Multiple concrete factories implement this interface.
  \item Each factory creates drivers for a specific platform or environment.
\end{itemize}
In our case:

\begin{narrowtable}{@{}ll@{}}
  \thead{Factory} & \thead{What it creates} \\
  \midrule
  \code{Esp32s3} & Real hardware drivers (for ESP32-S3) \\
  \code{Stub} & Simulated drivers for testing \\
\end{narrowtable}

\noindent The system logic only uses the factory interface and does not know which concrete factory
is used.

\subsection{Why do we switch to smart pointers?}
\label{c4:sec:whysmart}
In the previous example with raw pointers, the factory created objects using the \code{new}
operator:

\begin{cppcode}
return new gpio::Esp32s3(pin);
\end{cppcode}

\noindent The system logic had to delete the objects manually using the \code{delete} operator:

\begin{cppcode}
delete myLed;
\end{cppcode}

\noindent Using raw pointers introduces several problems:
\begin{itemize}
  \item It is easy to forget to call \code{delete}, which leads to memory leaks.
  \item Ownership of the object is unclear.
  \item Manual memory management makes programs harder to write correctly.
\end{itemize}
A modern C++ solution is to use \code{std::unique_ptr} from the \code{<memory>} header instead of
raw pointers.

A \code{std::unique_ptr}:
\begin{itemize}
  \item Represents exclusive ownership.
  \item Deletes the object automatically.
  \item Cannot be copied, only moved.
\end{itemize}
The \code{std::unique_ptr} returned by the factory transfers ownership of the created object to the
caller.

When the \code{std::unique_ptr} goes out of scope, the object is automatically destroyed.

If we use smart pointers, the factory can instead create objects by calling the
\code{std::make_unique()} function from \code{<memory>}:

\begin{cppcode}
return std::make_unique<gpio::Esp32s3>(pin);
\end{cppcode}

\noindent \code{std::make_unique()}:
\begin{itemize}
  \item Creates the object.
  \item Returns a \code{std::unique_ptr}.
  \item Ensures safe and clear construction.
\end{itemize}

\subsection{Exceptions and \code{noexcept}}
\label{c4:sec:noexcept}
In C++, memory allocation can fail. When we use \code{new} or \code{std::make_unique()}, an
exception (\code{std::bad_alloc}) may be thrown if the system cannot allocate memory.

In embedded systems, exceptions are often not used, or a predictable behavior is preferred when
errors occur. Therefore \code{noexcept} is still often used in driver and factory interfaces.

Example:

\begin{cppcode}
virtual std::unique_ptr<gpio::Interface> gpio(std::uint8_t pin) noexcept = 0;
\end{cppcode}

\noindent If a function marked with \code{noexcept} nevertheless throws an exception, the program
terminates immediately.

This provides clear and predictable error behavior. In many embedded systems this is desirable
because:
\begin{itemize}
  \item Running out of memory is often a critical error.
  \item The system cannot safely continue anyway.
\end{itemize}
For this reason, \code{noexcept} is used in our driver and factory interfaces. If a memory error
still occurs, the program will terminate, providing clear and predictable error behavior.

\subsection{Example implementation using Abstract Factory}
\label{c4:sec:smartexample}
Here we use a simplified form of the Abstract Factory pattern.

\subsubsection{Step 1: Factory interface}

\begin{cppcode}
#pragma once

#include <cstdint>
#include <memory>

namespace driver
{
/** GPIO driver interface. */
namespace gpio { class Interface; }
} // namespace driver

namespace driver::factory
{
class Interface
{
public:
    virtual ~Interface() noexcept = default;

    [[nodiscard]] virtual std::unique_ptr<gpio::Interface> gpio(std::uint8_t pin) noexcept = 0;
};
} // namespace driver::factory
\end{cppcode}

\noindent The factory now returns a \code{std::unique_ptr} instead of a raw pointer.

Note that \code{gpio()} is marked \code{[[nodiscard]]}: discarding the returned
\code{std::unique_ptr} would destroy the newly created object immediately, since nothing would be
left holding onto it. This is exactly the ``factory function returning an owning resource'' case
mentioned back in \secref{c1:sec:nodiscard}.

\subsubsection{Step 2: Real factory (for ESP32-S3)}

\begin{cppcode}
#pragma once

#include <memory>

#include "driver/factory/interface.hpp"
#include "driver/gpio/esp32s3.hpp"

namespace driver::factory
{
class Esp32s3 final : public Interface
{
public:
    Esp32s3() noexcept = default;
    ~Esp32s3() noexcept override = default;

    [[nodiscard]] std::unique_ptr<gpio::Interface> gpio(std::uint8_t pin) noexcept override
    {
        return std::make_unique<gpio::Esp32s3>(pin);
    }

    Esp32s3(const Esp32s3&)            = delete;
    Esp32s3(Esp32s3&&)                 = delete;
    Esp32s3& operator=(const Esp32s3&) = delete;
    Esp32s3& operator=(Esp32s3&&)      = delete;
};
} // namespace driver::factory
\end{cppcode}

\subsubsection{Step 3: Stub factory}

\begin{cppcode}
#pragma once

#include <cstdint>
#include <memory>

#include "driver/factory/interface.hpp"
#include "driver/gpio/stub.hpp"

namespace driver::factory
{
class Stub final : public Interface
{
public:
    Stub() noexcept = default;
    ~Stub() noexcept override = default;

    [[nodiscard]] std::unique_ptr<gpio::Interface> gpio(std::uint8_t pin) noexcept override
    {
        (void)(pin);
        return std::make_unique<gpio::Stub>();
    }

    Stub(const Stub&)            = delete;
    Stub(Stub&&)                 = delete;
    Stub& operator=(const Stub&) = delete;
    Stub& operator=(Stub&&)      = delete;
};
} // namespace driver::factory
\end{cppcode}

\subsubsection{Step 4: System logic with \code{unique_ptr}}

\begin{cppcode}
#pragma once

#include <cstdint>
#include <memory>

#include "driver/factory/interface.hpp"
#include "driver/gpio/interface.hpp"

namespace app::logic
{
class Logic final
{
public:
    explicit Logic(driver::factory::Interface& factory,
                   std::uint8_t ledPin,
                   std::uint8_t buttonPin) noexcept
        : myLed{factory.gpio(ledPin)}
        , myButton{factory.gpio(buttonPin)}
    {}

    void run() noexcept
    {
        bool buttonPrev{false};

        while (1)
        {
            // Read the current button input.
            const bool buttonCurrent{myButton->read()};

            // Toggle the LED on button press (rising edge).
            if (buttonCurrent && !buttonPrev)
            {
                myLed->toggle();
            }
            // Store the current button input for next comparison.
            buttonPrev = buttonCurrent;
        }
    }

    Logic()                        = delete; // No default constructor.
    Logic(const Logic&)            = delete; // No copy constructor.
    Logic(Logic&&)                 = delete; // No move constructor.
    Logic& operator=(const Logic&) = delete; // No copy assignment.
    Logic& operator=(Logic&&)      = delete; // No move assignment.

private:
    std::unique_ptr<driver::gpio::Interface> myLed;
    std::unique_ptr<driver::gpio::Interface> myButton;
};
} // namespace app::logic
\end{cppcode}

\note
\begin{itemize}
  \item No destructor is needed.
  \item The memory is released automatically.
\end{itemize}

\subsubsection{Step 5: \code{main}}
Just as in the factory example with raw pointers, we create a factory in the \code{main} function
and then pass it to the system logic to create objects.

Below is an example of how a factory can be used to run the system logic on ESP32-S3 by passing an
instance of \code{driver::factory::Esp32s3} to the system logic and then calling \code{run()}:

\begin{cppcode}
#include <cstdint>

#include "driver/factory/esp32s3.hpp"
#include "app/logic/logic.hpp"

int main()
{
    constexpr std::uint8_t ledPin{2U};
    constexpr std::uint8_t buttonPin{3U};

    // Create system logic and initialize the system.
    driver::factory::Esp32s3 factory{};
    app::logic::Logic logic{factory, ledPin, buttonPin};

    // Run the system continuously.
    logic.run();
    return 0;
}
\end{cppcode}

\noindent As shown in the example factory with raw pointers, it is sufficient to switch to the stub
factory \code{driver::factory::Stub} in order to run with stub drivers:

\begin{cppcode}
#include <cstdint>

#include "driver/factory/stub.hpp"
#include "app/logic/logic.hpp"

int main()
{
    constexpr std::uint8_t ledPin{2U};
    constexpr std::uint8_t buttonPin{3U};

    // Create system logic and initialize the system.
    driver::factory::Stub factory{};
    app::logic::Logic logic{factory, ledPin, buttonPin};

    // Run the system continuously.
    logic.run();
    return 0;
}
\end{cppcode}

\subsection{Summary}
\label{c4:sec:smartsummary}
In this version we use:
\begin{itemize}
  \item \textbf{Interfaces} for all drivers.
  \item \textbf{Abstract Factory} to create the correct type of drivers.
  \item \textbf{\code{std::unique_ptr}} for safe ownership.
\end{itemize}
Advantages:
\begin{itemize}
  \item No manual \code{delete}.
  \item No memory leaks.
  \item Clear ownership.
  \item The same system logic works:
  \begin{itemize}
    \item On real hardware.
    \item In simulation.
  \end{itemize}
\end{itemize}
